% Using InitechBase, Chapter 2: Using Commands (annex Sec 8.2 U2).  Sources:
% samir_main.c (the REPL line, 254 characters, verb splitting, USE/CLOSE/DO
% file), cmd/flow.c (verb dispatch, STORE/=, memvars), core/lex.c and
% core/parse.c (constants, operators, precedence), core/eval.c (operator
% semantics, the coercion table spec/samir/xbase_coercion.json),
% cmd/query.c and cmd/mutate.c (scope, FOR, WHILE), cmd/workarea.c and
% cmd/nav.c (work areas, record pointer, BOF/EOF), fs/ndx.c (indexes).
\chapter{Using Commands}
\Gates{test-pubs-examples, test-xbase-lex, test-xbase-parse, test-xbase-eval, test-interp-nav}
This chapter describes how InitechBase commands are written: the parts of a
command, the expressions that commands work with, memory variables, work
areas, the record pointer and index files.  Chapter~5 gives the form of each
command.

\section{Typing a Command}
A command is typed at the dot prompt on one line of up to 254 characters and
ends when you press \key{Enter}.  It begins with the command word, the
\term{verb}, which says what is to be done: \cmd{LIST}, \cmd{USE},
\cmd{REPLACE}.  Some commands consist of the verb alone; most continue with
the things the verb works on and with clauses that limit it.

Command words, field names and memory variable names may be typed in capital
or small letters.  InitechBase does not distinguish them, but it keeps the
expressions of a \cmd{LIST} heading as you typed them:
\DotExample{examples/u2-case}
Type command words in full.  Separate the parts of a command with blanks;
more than one blank does no harm.  To correct a mistake before you press
\key{Enter}, press \key{Backspace}; Appendix~A lists the keys you can use.

A command that InitechBase does not know is refused with
\xmsg{16}{*** Unrecognized command verb.}  A command it knows but cannot carry
out is refused with a message that names the problem (Chapter~7).  In both
cases nothing has been changed, and the dot prompt is displayed again.

\section{The Parts of a Command}
The full form of a command that processes records is:

\SyntaxLine{\xp{verb} \opt{\xp{scope}} \opt{\xp{expression list}}
  \opt{\kw{FOR} \xp{condition}} \opt{\kw{WHILE} \xp{condition}}}

The \term{scope} names the records the command works on.  The
\term{expression list} names the fields or the values it shows or uses.  The
\kw{FOR} and \kw{WHILE} clauses give the conditions that the records must
meet.  Each part is described below.  The scope comes directly after the verb;
\kw{FOR} and \kw{WHILE} follow the expression list.

\section{Expressions}
An \term{expression} is a value, or a combination of values that InitechBase
calculates.  It may be a constant, a field name, a memory variable, a
function, or any of these joined by operators.  Every expression has one of
four types: character, numeric, date or logical.  A memo field is used as a
character value.

\SideHead{Constants}A character constant is enclosed in single quotation marks,
double quotation marks or square brackets; the delimiters are not part of the
value.  A numeric constant is written with digits, a decimal point and a minus
sign as needed.  The logical constants are \texttt{.T.} and \texttt{.F.};
\texttt{.Y.} and \texttt{.N.} may be written for them.  A date is written as a
character string converted by \fn{CTOD()}.
\DotExample{examples/u2-constants}

\SideHead{Fields}The name of a field stands for the value of that field in the
current record of the database in the selected work area.

\SideHead{Memory variables}The name of a memory variable stands for its value
(see ``Memory Variables'' below).

\SideHead{Functions}A function computes a value from its arguments; Chapter~6
describes each one.

\section{Operators}
\SideHead{Arithmetic}On numbers: \texttt{+} addition, \texttt{-} subtraction,
\texttt{*} multiplication, \texttt{/} division, \texttt{**} or \texttt{\^{}}
exponentiation.  Exponentiation is done first, then multiplication and
division, then addition and subtraction; operators of the same rank are done
from left to right.  Parentheses change the order.
\DotExample{examples/u2-operators}
The result of a calculation is shown with as many decimal places as SET
DECIMALS gives, two when InitechBase starts; a result that is a whole number
is shown without decimal places.  The value itself keeps its full precision:
\DotExample{examples/u2-decimals}

\SideHead{Character}\texttt{+} joins two strings as they are.  \texttt{-}
joins them and moves the trailing blanks of the first to the end of the
result.  \texttt{\$} is true when the first string occurs in the second;
capital and small letters are different.
\DotExample{examples/u2-strings}

\SideHead{Date}A number of days may be added to a date or subtracted from it;
the result is a date.  One date subtracted from another gives the number of
days between them.
\DotExample{examples/u2-dates}
Date arithmetic follows the calendar, leap years included:
\DotExample{examples/u2-datecalc}

\SideHead{Relational}\texttt{<} less than, \texttt{>} greater than,
\texttt{=} equal, \texttt{<>} or \texttt{\#} not equal, \texttt{<=} less than or
equal, \texttt{>=} greater than or equal.  Both sides must be of the same type;
the result is logical.  Characters are compared by their codes (Appendix~D),
so capital letters come before small letters.  For \texttt{=} between
character strings, see SET EXACT in Chapter~5.  While SET EXACT is OFF, \texttt{=}
is true when the string on the left begins with the string on the right, so
the order of the two sides matters, and an empty string on the right equals
any string.  The characters are compared one by one, so \texttt{'10'} is less
than \texttt{'9'}:
\DotExample{examples/u2-compare}

\SideHead{Logical}\texttt{.AND.}, \texttt{.OR.} and \texttt{.NOT.} combine
logical values.  \texttt{.NOT.} is done first, then \texttt{.AND.}, then
\texttt{.OR.}.
\DotExample{examples/u2-logical}

Operands of the wrong type are refused with \xmsg{9}{Data type mismatch.}
InitechBase does not convert numbers to character strings or the other way
round: use \fn{STR()}, \fn{VAL()}, \fn{DTOC()} and \fn{CTOD()}.  Appendix~E
lists every combination of types for every operator.  A division by zero
gives \xmsg{39}{Numeric overflow (data was lost).}

\section{The Scope}
The scope names the records a command works on.

\SideHead{ALL}Every record of the database, beginning with the first.

\SideHead{NEXT \xp{n}}\xp{n} records, beginning with the current record.

\SideHead{RECORD \xp{n}}Only the record with record number \xp{n}.

\SideHead{REST}The current record and every record after it.

Each command has a default scope, used when no scope is given: \cmd{LIST},
\cmd{LOCATE} and \cmd{REPORT FORM} work on all the records; \cmd{DISPLAY},
\cmd{REPLACE}, \cmd{DELETE} and \cmd{RECALL} on the current record.  The
commands that change records take the scope \kw{ALL} when a \kw{FOR} clause is
given without a scope.
\DotExample{examples/u2-scope}

\section{Conditions}
A \term{condition} is a logical expression tested for each record.  InitechBase
does not treat a number or a string as true or false: a condition of another
type is refused with \xmsg{37}{Not a Logical expression.}

\SideHead{FOR}A command with \kw{FOR} \xp{condition} examines every record of
its scope and processes those for which the condition is true.

\SideHead{WHILE}A command with \kw{WHILE} \xp{condition} processes records, in
order, as long as the condition is true, and stops at the first record for
which it is false.  Begin at the first record wanted and give the scope
\kw{REST}.  With an index on the field tested, \kw{WHILE} processes a group
of records without reading the rest of the database:
\DotExample{examples/u2-while}

\section{Memory Variables}
A \term{memory variable} holds a value that is not part of any database: a
total, an answer typed by the user, a date.  It is created when a value is
first stored in it with \cmd{STORE} or with \texttt{=}, and it takes the type
of that value.  A memory variable name has up to 10 characters: letters,
digits and underscores, beginning with a letter.  InitechBase holds up to 256
memory variables at once.
\DotExample{examples/u2-memvar}
A memory variable exists until it is released with \cmd{RELEASE} or until the
procedure that created it returns.  \cmd{PUBLIC} and \cmd{PRIVATE} decide how
long memory variables live in programs (Chapter~3).

When a field and a memory variable have the same name, the name stands for the
field.  Give memory variables names that are not the names of fields: the
programs of this manual begin them with the letter M.

\section{Work Areas}
InitechBase can keep ten databases open at once, each in its own
\term{work area}, numbered 1 to 10.  \cmd{SELECT} chooses the work area in
which the next commands work; \cmd{USE} opens a database in the selected work
area.  The database of a work area is known by its \term{alias}: its file name
without the extension, or the name given with \kw{ALIAS}.

Each work area keeps its own record pointer and its own index files.  To use a
value from one database while working in another, store it in a memory
variable first:
\DotExample{examples/u2-areas}

\section{The Record Pointer}
In each work area, InitechBase keeps a \term{record pointer} on one record, the
\term{current record}.  When a database is opened, the pointer is on the first
record.  \cmd{GO}, \cmd{SKIP}, \cmd{LOCATE}, \cmd{CONTINUE}, \cmd{SEEK} and
\cmd{FIND} move it; \cmd{LIST} and the other commands that process a scope
leave it at the end of the scope.  \fn{RECNO()} gives the number of the
current record.

After the last record there is one more position, the \term{end of file}.
When the pointer is there, \fn{EOF()} is true and \fn{RECNO()} is one more
than the number of records.  Before the first record is the \term{beginning of
file}: a \cmd{SKIP} back from the first record makes \fn{BOF()} true.
Appendix~B lists what each command does at the beginning and end of the file.
\DotExample{examples/u2-pointer}

\section{Index Files}
An \term{index file} holds the records of a database in the order of a
\term{key}, such as the last name, the department or the salary.  When a
database is opened with an index file, \cmd{LIST}, \cmd{SKIP}, \cmd{GO TOP} and
the other commands take the records in key order, and \cmd{SEEK} and \cmd{FIND}
find a record by its key at once.  The record numbers do not change.
\DotExample{examples/u2-index}
Up to seven index files may be open with one database.  The first named in
\cmd{USE} is the \term{master index}, which gives the order; SET ORDER chooses
another.  While index files are open, InitechBase keeps them up to date when
records are added with \cmd{APPEND BLANK} or changed with \cmd{REPLACE}, and
rebuilds them when \cmd{PACK} removes records.

The index files of the sample diskette are:

\noindent\begin{tabular}{@{}lll@{}}
\hline
\textsf{\textbf{Index file}} & \textsf{\textbf{Database}} & \textsf{\textbf{Key}}\\
\hline
\file{STAFFNAM.NDX} & \file{STAFF.DBF} & \texttt{LASTNAME}\\
\file{STAFFNO.NDX} & \file{STAFF.DBF} & \texttt{EMPNO}\\
\file{STAFFDPT.NDX} & \file{STAFF.DBF} & \texttt{DEPT}\\
\file{STAFFSAL.NDX} & \file{STAFF.DBF} & \texttt{SALARY}\\
\file{DEPTS.NDX} & \file{DEPTS.DBF} & \texttt{DEPT}\\
\file{INVNO.NDX} & \file{INVOICE.DBF} & \texttt{INVNO}\\
\file{INVCUST.NDX} & \file{INVOICE.DBF} & \texttt{CUST}\\
\hline
\end{tabular}

An index file belongs to one database: open it only with that database.
